% =====================================================================
%  R. Nielsen, PERSONLIGE GRUNDE
%  In: For Ide og Virkelighed. Et Tidsskrift, udgivet af R. Nielsen, Bj. Bjørnson og Rud. Schmidt.  Aargang 1869.
%  Kjøbenhavn: Chr. Steen & Søns Forlag; Louis Kleins Bogtrykkeri.  pp. 585-622.  Antiqua.
%
%  TRANSCRIPTION -- GENERATED; DO NOT EDIT BY HAND.  The five Nielsen essays of the 1869 volume are built by ONE pipeline
%  (ebook-method/for-idee-1869/) into a combined working file, which split5.py cuts into the five transcriptions.
%  Rebuild from the repo root:   sh ebook-method/for-idee-1869/rebuild.sh
%
%  STATUS: COMPLETE (2026-10-13): the volume's title page and the whole essay, pp. 585-622, with its notes.
%
%  EDITION.  The only printing.  Text of record: the Internet Archive's scan of the University of Toronto copy
%  (forideogvirkelig00kjobuoft), ~/bibliotek/Nielsen, Rasmus/1869-for-idee-og-virkelighed.pdf (SCANS.tsv), 886 PDF pages,
%  400 ppi colour; text layer = the IA's ABBYY FineReader OCR.  Printed page N = PDF N + 10 to p. 688; a plate is bound in
%  before p. 689, so pp. 689-706 = PDF N + 12.  Title page PDF 7.  No e-book, no other scan, no earlier transcription found
%  (checked 2026-10-13: the repo, the web, the Internet Archive; HathiTrust and Google Books could not be queried that day).
%
%  ACCURACY (2026-10-13; one collation for all five essays, 189 pp., ~48,000 words).  The words are those of the ABBYY layer
%  (draft.py), collated word by word with an independent tesseract reading of the images (collate/check.py, witness 'ocr'):
%  1132 disagreements at the first collation.  check.py's own rules settled 96; rules.py settled ~330 more, each rule
%  switched on only after a BLIND PILOT of 107 crops (reader A, 21/22 controls) had measured it on this book: tesseract's
%  glyph faults 8/8, aa read an 8/8, accents it adds 8/8, the layer's glyph faults 7/7 and 8/8, o/ø by the ø rule 8/8, the
%  printer's space before a mark 8/8 (house rule), and the narrow word-division rules; the quote, dash, note-number and
%  word-seam rules failed the pilot and stayed off.  The rest were read at the image BLIND (sheets.py: 579 word crops on
%  58 sheets, two planted controls a sheet, 116 of 116 controls right; six readers) and about 110 spots in the main
%  session from line crops (zoom.py).  STRUCTURE read at the image (reader K): the essays' headings, the volume's title
%  page, every footnote whole (the notes were collated where they stand and then moved to their marks), the signatures.
%  LETTERSPACING (the print's emphasis, \emph): candidate lines by box width (spacedwidth.py; the gap measure is noise on
%  this scan): 163 lines read blind (22/22 roman controls), doubtful ones read in the main session.  ITALICS by stroke
%  slant (check.py --italics): 24 crops read blind; the only italics are Latin quotations.
%  SWEEPS after OPEN 0, each spot read at the image: the W-shortens / T-longer list of the blind decisions (found two
%  W decisions that had dropped a note's mark and a closing «, and one that had joined „der f.“); --moved and --splits;
%  the split sweep; the reader-line sweep; paragraph starts against tesseract's indents; quotation marks (the print
%  closes with « and often with », both as printed); the unknown-word sweep (31 spots: 16 corrections).
%  TRANSLATION as a last reading: the twelve translators listed 79 odd spots (ebook-method/for-idee-1869/tr/odd-*.tsv),
%  ~60 read at the image (zoom zo, zw, zx): mended -- f's whose head did not ink, read t by both OCRs (pp. 21, 36, 309, 319,
%  441; a recurring defective sort, as in Folkelige Foredrag), two blind decisions that had taken a reader's t for the
%  print's f (pp. 322, 539), a speck read as a stop (p. 589), a lost dash (p. 525), the verse on p. 414 set as one quotation,
%  false paragraph breaks at page turns left by the notes' removal (pp. 430, 553, 620, 695, 703); the rest are the print's.
%  SAMPLED PAGE CHECK (pagecheck.py, seed 202610136): pp. 12, 418, 521, 694 and the title page read whole and blind, 1069
%  words: 0 word and 0 punctuation errors; one letterspaced passage in a note missed (p. 694), mended.
%  What this cannot see: a misprint both OCRs normalise into a real word.
%  The ø rule (Hans 2026-10-08): o-forms that are no word (Oieblik, Ifolge, imode ...) are set with ø.
%  PRINTED AS IS p. 605: „højærværdige“ with j
%  PRINTED AS IS p. 607: „hundede“ for hundrede
%  PRINTED AS IS p. 618: „aandige“
%  PRINTED AS IS p. 586: „ogaa“ for ogsaa
%  DEFECT, letter supplied p. 619: „indrømme“: the i did not ink (a speck); letter supplied
%  Not transcribed: running heads, folios, sheet signatures („For Ide og Virkelighed.  20“), the library's stamp; the large
%  initial of each essay is set as a plain letter.  Quotation marks as printed (» « in the text, „ “ in the notes).  The notes
%  are numbered ¹) per page, as printed (p. 553: *)).
%
%  TOOLS (ebook-method/for-idee-1869/, copied from folkelige-foredrag): draft.py, ocr.py, rules.py, sheets.py + reads.py +
%  READER-BRIEF.md, heads.py, spacedwidth.py + spacedlines.py + spaced.tsv, zoom.py, pagecheck.py + PAGECHECK-BRIEF.md,
%  patches.py (PATCHES, NOTES, POST), decisions.tsv, the sweeps, split5.py, rebuild.sh.
% =====================================================================
\documentclass[12pt,a4paper,openany]{book}
\usepackage[utf8]{inputenc}
\usepackage[T1]{fontenc}
\usepackage{libertinus}
\usepackage{libertinust1math}
\usepackage[danish]{babel}
\usepackage[protrusion=true,expansion=true]{microtype}
\usepackage{setspace}
\usepackage[margin=1.3in]{geometry}
\usepackage{fancyhdr}
\usepackage{hyperref}
\hypersetup{pdftitle={Personlige Grunde}, pdfauthor={Rasmus Nielsen}, hidelinks}
\pagestyle{fancy}
\fancyhf{}
\fancyhead[LE,RO]{\thepage}
\fancyhead[C]{\textit{R. Nielsen: Personlige Grunde (1869)}}
\setstretch{1.3}
\usepackage{marginnote}
\renewcommand{\thefootnote}{\arabic{footnote})}   % the print numbers each page's note ¹) (p. 553: *))
\newcommand{\opage}[1]{\marginnote{\footnotesize\textit{[#1]}}}
% the essay's heading: title (and subtitle), centred, a short rule; the print's large initial is not reproduced
\newcommand{\essayhead}[2]{\begin{center}{\large #1}\ifx&#2&\else\\[0.3em]{\normalsize #2}\fi\\[0.5em]\rule{3em}{0.4pt}\end{center}\medskip\noindent}
\newcommand{\signature}[1]{\par\smallskip\hfill #1\hspace*{2em}\par}

\begin{document}

% --- front matter: the volume's title page (PDF 7), read at the image (reader K, F7; the page check reads it again blind).
\begin{titlepage}
\begin{center}
\vspace*{3em}
{\Huge For Ide og Virkelighed.}\\[1em]
\rule{8em}{0.4pt}\\[1em]
{\large Et Tidsskrift,}\\[0.8em]
{\small udgivet}\\[0.4em]
{\footnotesize af}\\[0.6em]
{\large R. Nielsen,}\\[0.4em]
{\large Bj. Bjørnson og Rud. Schmidt.}\\[1em]
\rule{8em}{0.4pt}\\[1em]
{\large Aargang 1869.}\\[1em]
\rule{8em}{0.4pt}\\[4em]
Kjøbenhavn.\\[0.4em]
Chr. Steen \& Søns Forlag.\\[1.5em]
{\small Louis Kleins Bogtrykkeri.}
\end{center}
\end{titlepage}

% ===== TEXT (pp. 585-622) =====
\essayhead{Personlige Grunde.}{}
% --- p. 585 ---
\setcounter{footnote}{0}%
\opage{585}Enhver, der kjender lidt til sig selv og Verden, vil
vistnok have gjort den Erfaring, at der gives en Overbevisning,
der vanskelig lader sig fremsætte og retfærdiggjøre
for Andre, og det ikke af Mangel paa Grunde,
nei tvertimod, af Overflødighed paa Grunde. Det gjælder
i de fleste Tilfælde om kvindelig Overbevisning. En
ung Pige, f. Ex., kan om et høist vigtigt Anliggende
være ganske paa det Rene med sig selv, fuldkommen
vis i sin Sag, og dog, naar hun skal forklare sig, er
det, som sank hun i en Strøm af Grunde, hun forstummer.
Hun kan ikke udelukkende beraabe sig paa nogen
enkelt Grund, og vilde hun forsøge paa at opregne alle
sine Grunde, da kunde hun ligesaa godt, det føler
hun, begynde at tælle Bladene i Skoven eller Draaberne
i Havet. At hendes Grunde ere personlige, ikke
videnskabelige, er vel en overflødig Bemærkning. Man
vilde imidlertid tage storligen Feil, om man bildte sig
ind, at en Overbevisning af subjektive, personlige Grunde
kun egnede sig for Kvinder, ikke for Mænd. Absolut
personlige Grunde gjælde overalt, hvor der handles om
en absolut personlig Sag.

I Øhlenschlægers Axel og Valborg er det, som
bekjendt, Sortebroder Knud, der i den Omstændighed,
% --- p. 586 ---
\setcounter{footnote}{0}%
\opage{586}at Brud og Brudgom ere »Daabsøskinde«, skarpsindig
opdager en efter Kirkens Ret fuldgyldig Grund til
at hindre Vielsen; at Grunden er objektiv, en Grund,
der ikke tilsteder Indsigelser, faaer Axel at vide, da
han, der »ogsaa kan læse«, seer med egne Øine, at det
staaer i Kirkebogen. Den objektive Grund overbeviser
ikke Axel om, at hans Kjærlighed til Valborg er syndig
--- i saa Fald maatte Forbudet være opfattet fra den
religiøse Side, altsaa subjektivt --- nei den overbeviser ham
ikke, men den viser ham en fjendlig Modstand, en Skranke,
haard som Skæbnen, ubøielig som Nødvendigheden; den
viser ham Umuligheden af at naae sin Længsels Maal, sit
Eden, »sit salige Valhalla«. Saaledes kan en eneste objektiv
Grund trykke som Dødens Haand paa den Levendes
Bryst og kue alle Hjertets Grunde, om disse end ere
nok saa inderlige, nok saa stærke. Det Tragiske i Katastrofen
er et Udtryk for, at subjektive og objektive
Grunde ere uensartede.

Men et saadant Udtryk kan ogaa faaes paa en anden   % PRINTED AS IS (p. 586): „ogaa“ for ogsaa
Maade. Vi kunne jo forestille os, at Axels Ven Vilhelm
ved et af de mange vidunderlige Tilfælde, der altid
staae til Digterfantasiens Raadighed, havde i rette Tid
faaet Nys om Hemmeligheden, var kommen til Rom
faa Dage efter Axels Afreise og havde saa i al Stilhed
virket for, at den oversete Hindring blev hævet.
Just som Knud har aabnet Kirkebogen, og Axel styrter
for Valborgs Fødder med Fortvivlelsens Udraab:
»Valborg! Valborg! det er forbi«, træder da Vilhelm op,
griber Valborgs Replik: Nei, nei! og fortsætter: Hadrianus
Episcopus, servus servorum Dei,\dots{} kort at melde:
Hans Hellighed har af synderlig Naade tilladt, at Axel
% --- p. 587 ---
\setcounter{footnote}{0}%
\opage{587}Thordsøn ægter Valborg Immers Datter, uanseet, at de
ere Gudsødskende, dog paa den Betingelse, at samme
Axel Thordsøn først med klare, uigjendrivelige Grunde
viser, hvorfor han netop elsker Valborg og aldrig vil
kunne elske nogen anden Kvinde. --- Budskabet er
en glædelig Overraskelse for det nylig saa ulykkelige
Par; Axel spørger i Henrykkelse: Er det det Hele,
forlanger Paven intet Andet? Vilhelm. Nei, intet Andet.
Og dog er Sagen besynderlig nok. Det var en
Kardinal, der gik i Forbøn for Dig, et fint Hoved, en
subtil Skolastiker. Nu maa Du vide, hvad Du vist
ikke har lagt Mærke til, at der iblandt de Lærde føres
en heftig Strid om Villiens Forhold til Bevæggrundene,
om den menneskelige Friheds Forhold til Nødvendigheden,
idet Nogle paastaae, at den sande Frihed i Grunden
ikke er Andet end en Slags sjælelig Naturnødvendighed,
medens Andre mene, at der er en dybtgaaende
Forskjel mellem Nødvendighed og Frihed, at den frie
Villie kun bestemmes af Grunde, forsaavidt den bestemmer
sig selv. Kardinalen forestillede nu Paven, at det
meget vilde bidrage til Sagens Opklaring, om nogle
iblandt de mange Elskere, der aarlig overhænge den
hellige Fader med Bøn om Dispensation, bleve opfordrede
til at udtale sig nærmere om deres Grunde. Thi Elskere
tale ofte et tvetydigt Sprog, naar de tale om Grunde.
Snart sige de, at de elske af Nødvendighed, at de ikke
kunne Andet; det er da Livet om at gjøre, og man maa
see at hjælpe paa de Ulykkelige, saa godt man kan.
Men saa sige de igjen, at de elske med Frihed, at Kjærlighed
er en fri Sag, at de med fuld Beslutning have
valgt den Udvalgte, kaaret den Udkaarne, at de ikke
% --- p. 588 ---
\setcounter{footnote}{0}%
\opage{588}kunne Andet, fordi de ikke \emph{ville} Andet; men naar Valget
saaledes afhænger af Villien og skeer med Frihed,
saa kunde Supplikanten jo ogsaa vælge en Anden, og der
er ingen paatrængende Grund til at give ham Dispensation.
Ved at samle og sammenligne de mange Udtalelser kunde
Hans Hellighed da erfare om Villien i dette vigtige Punkt
bevæges af »det Dunkle, det Egoistiske og Driftagtige«,
der »hænger sammen med den individuelle Naturværens
uoplukkede Enhed«, eller af »Videns Almene, Væsensuniversalitetens
Bestemmelse, der udtrykker sig i den
Forsigværen i den frigjorte Almindeligheds Form, der er
Selvbevidsthedens«, med andre Ord, om de virkelig elske
af Nødvendighed eller af Frihed, eller tildels af Nødvendighed
tildels af Frihed, eller\dots{} Axel. Hold
dog op med dette lange Slæng af lærde Fraser; Du
veed, jeg er ingen Klerk. Jeg gaaer til Sagen ligefrem;
jeg har en Grund, en eneste enfoldig; min Grund
vil Paven ikke forkaste, det er min Valborg selv ---
»granvoxne Valborg, Rose for alle Blomme«! Vilhelm.
Ja min Ven, det er en lyrisk Grund; den Grund kan Du
ikke byde hans Hellighed; det var heller ikke den Slags
Grunde, hvorom Kardinalen talede, han mente nemlig
objektive, videnskabelige Grunde. Axel. Saa maa Du
hjælpe mig at mynte Grunden med det rette Præg; thi
jeg vil sige Dig, det er med min Grund som med Draupner,
Odins Ring: den er kun een og bliver dog til mange.
Tager Paven ikke imod den ene Grund, saa maa han
tage imod den anden. Jeg har mange Grunde, Vilhelm,
naar jeg tænker efter, mange; »Havet har ikke saa mange
Dyr i sine Bølger, aldrig saae Natten saa mange Stjerner
paa Himlens Hvælving, ikke boer der saa mange
% --- p. 589 ---
\setcounter{footnote}{0}%
\opage{589}Fugle i Skoven, ikke saa mange Straa har Mark og
Eng«, som mit Hjerte har Grunde, hvi jeg elsker Valborg.
--- Saa begynder Overveielsen. Vielsen er udsat, Axel
grunder, Valborg venter --- hun faaer at vente længe.

Saameget er vel heraf klart, at en subjektiv personlig
Grund aldrig kan blive upersonlig, saa lidt som
omvendt. Kunde de uensartede Grunde i videnskabelig
Forstand sammensmelte til Eet, da vilde den tænkende
Bevidsthed ogsaa være i Stand til at omsætte
den ene i den anden, forvandle den ene ved Hjælp af
den anden. Men hvor umulig en saadan Forvandling er,
kan sees af de anførte Exempler. Den objektive Grund,
at Daabssødskende, ifølge Kirkens Forbud, ikke tør ægte
hinanden, kan, naar den virker i al sin Objektivitet, i
det Høieste kun bevirke, at de Elskende bøie sig under
Tilskikkelsen, at de i stille Taalmod give Afkald paa
Livets Lykke. Det er en Ydmyghedens Kraft i Kvindens
Sjæl, at hun inderligere formaaer at bøie sig under
det Uundgaaelige, end Manden; vi høre Valborg tage
Ordet for den stille Resignation. »Stille skal jeg sidde
saadan i en lille Celle, virke Guld i Silke, henleve blid
og sørgelig min Tid, som Turtelduen«... Som hun
trøster sig selv, vil hun ogsaa trøste sin Ven, at han
skal være taalmodig, stille. »Du skal drage til din
Gaard, op til din Søster, den gode Helfred. Tiden læger
i Verden Alt, den læge dine Saar. Naturen i sin ædle,
stille Storhed, en kjærlig Søsters Omgang lindre Hjertet
mer end det muntre Selskab«. Det Subjektive er herved
ingenlunde blevet objektivt, men den personlige Grund
er ved Sammenstød med Skæbnen forvandlet i sig selv.

Det er jo ikke Valborgs Mening, naar hun taler om,
% --- p. 590 ---
\setcounter{footnote}{0}%
\opage{590}at Tiden læger alle Saar, at hun med Tiden vil glemme
Axel, eller at han skulde glemme hende. Hvad der
forandrer sig med Tiden er egenlig kun det Udvortes,
ikke det Indvortes. Udvortes, for Verdens Øine maa
hun skilles fra den Elskede, men hun taber ham ikke;
hun kjender et Sted, hvor han altid er, et fredhelligt
Sted i det Indre. Derinde, hvor han altid er, der
bygger hun Kirken --- med gamle Grave for kjære
Minder --- der lever hun, der vaager hun og holder
daglig Gudstjeneste: krandser Søilen, hilser Pilgrimmen,
græder ved hans Bryst, følger ham til Alteret,
til Sorgens Alter, til Vielsen, den Vielse, som
ingen Sorteknud afbryder. Det er jo kun subjektivt,
Altsammen subjektivt, men det er personligt. Skulde
det Subjektive derimod forvandles til det Objektive, det
Personlige til det Upersonlige: da maatte Sagen tage
en anden Vending. Det maatte da være, at Axel pludselig
ved Skæbnens Stød var bleven objektiv Tænker,
var kommen til at indsee, at det Oprindelige i Menneskets
personlige Væren ikke er Hjerte og Villie, men
Tænken og Viden; han maatte da forklare sig for Valborg.
Forstaa mig Valborg, det haarde Slag, der rammer
os, have vi selv forskyldt; vi have begaaet en stor
Synd, ikke mod Gud i Himlen --- det er en subjektiv
Sag --- men mod Tænkens og Værens evige Love, mod
Metafysiken. Seer Du, vi have i vor Ungdommelighed
begyndt med vor egen Væren, med »Væren i Individualitetens
Form«, vi have anseet denne vor indskrænkede
Subjektivitets Tilfældighed som et tilforladeligt Udtryk
for »individuel Væren i Almindelighed«, vi have desværre
taget Existensens Begreb i »en vilkaarlig Dob%
% --- p. 591 ---
\setcounter{footnote}{0}%
\opage{591}beltbetydning«: det er en Uklarhed i Tanken, »en Tilsnigelse«,
der hævner sig, nu da vi vaagne til Bevidsthed.
Forstaa mig: »al Væren i Individualitetens Form
er Existens; der kan altsaa gives en saadan Existens,
der ikke er Bevidsthed, ikke theoretisk. Men naar \emph{Bevidsthed}
bliver Existensens Subjekt, saa indeslutter
Existens Theorien, »thi Bevidstheden er kun existerende
som værende theoretisk«. --- Du stirrer paa mig, Du
forstaaer mig ikke; giv Tid! Du skal vel komme til at
forstaae mig. At en Yngling elsker en Pige, kan forsvares
af objektive Grunde; det er en Manifestation af
Naturens Almene; det er Slægtens Liv, der fortsætter
sig igjennem Individerne, hvis Existens forsaavidt maa
siges at gaae ind under individuel Væren i Almindelighed;
men at to Forelskede skulde have Ret til at forgude
hinanden i den Grad, at hun fortvivler og vil gaae
i Kloster, hvis hun ikke kan faae ham, just ham,
uagtet der er andre unge Mænd med ligesaa elskværdige
Egenskaber, der tilbyde sig; og at han bliver desperat
og vil i Krig og absolut slaaes ihjel, naar han ikke kan
faae hende, just hende, den Udvalgte: seer Du, Valborg,
en saadan Udvælgelse lader sig ikke forsvare af objektive
Grunde, det er at sætte »Væren i Individualitetens
Form« for »individuel Væren i Almindelighed«, det er
at krænke »Theorien«. Viis nu i Gjerningen, at Du har
begrebet det Objektive, at din Bevidsthed omsider er
bleven din Existenses Subjekt«, og at dette dit Subjekt
virkelig indeslutter Theorien: saa bliver Du Hakons Brud
--- til Ære for det Objektive! Norges Dronning --- til
Held og Hæder for det Almene! Saa er jeg trøstet;
Du vil ogsaa være det. Farvel!

% --- p. 592 ---
\setcounter{footnote}{0}%
\opage{592}Hvorledes Valborg vil opfatte Forklaringen, om hun
vil antage sig selv eller den Talende for afsindig, om
hun i brusende Vrede vil reise sig mod den Troløse, eller
om hun svimlende vil synke sammen i Afmagt, derom kan
der være forskjellig Mening; men saameget er vist: hun
vil ikke forstaae Forklaringen, ikke fatte Theoriens Objektive.
Hun gav ham sit Hjerte, sin Tro; dette Ene forstaaer
hun, dette Subjektive; men det Objektive? hvor
skulde hun forstaae »Theoriens Objektive«: hun er jo
--- kun en Kvinde!

»Kjærligheden«, siger S. Kierkegaard, »har dog i
Digterne sine Præster, og stundom hører man en Røst,
der veed at holde den i Hævd; men om Troen høres
der intet Ord, hvo taler til denne Lidenskabs Ære? Filosofien
gaaer videre. Theologien sidder sminket ved
Vinduet og beiler til dens Gunst, falbyder sin Deilighed
til Filosofien. Det skal være vanskeligt at forstaae
Hegel, det er et Underværk, men at komme ud over
Abraham, det er en smal Sag. At gaae ud over Hegel,
det er et Underværk, men at komme ud over Abraham
det er det Letteste af Alt.« Om nu Kierkegaard just
har Ret, naar han kalder Troen en Lidenskab, er et
Spørgsmaal for sig; men der er Et, hvori han ubetinget
har Ret, dette nemlig, at Troen ligesaa vel som
Kjærligheden er en personlig Sag. At ville i dette
Kapitel overhale en saadan Forfatter med Skolekritik
paa et Par Blade og give det Udseende af, at man
gaaer videre, gaaer ud over Kierkegaard, medens
man i Virkeligheden ikke blot gaaer tilbage til Hegel,
men endogsaa bagved Hegel, er ikke til Sandhedens
Fremme. Det er ikke for Intet, at Kierkegaard i alle
% --- p. 593 ---
\setcounter{footnote}{0}%
\opage{593}sine Skrifter, ikke blot i de æsthetiske og filosofiske,
men ogsaa i de religiøse, opbyggelige, viser hen til
det Erotiske og tager Exempler af naturlige Kjærlighedsforhold.
Han vil herved gjøre det ret iøinefaldende,
hvor skarp Modsætningen mellem det Personlige
og det Upersonlige er. Thi hvad der for Liv og Personlighed
er det Vigtigste, det Allervæsenligste: \emph{dette}
Menneske, \emph{denne} Sjæl, dette Hjerte, maa for den Bevidsthed,
»der kun existerer som værende theoretisk«
nødvendigvis blive til noget aldeles Uvæsenligt, noget
Forsvindende og desaarsag Ubetydeligt. Theorien naaer
kun til \emph{en} Sjæl, aldrig til \emph{denne} Sjæl, kun til \emph{et} Hjerte
i Almindelighed, Hjertets Begreb, aldrig til \emph{dette} Hjerte.
Hvad Hjertet i og for sig, det objektive Hjerte, attraaer
og i Kraft af Begrebet maa attraae, er Hovedsagen; hvad
et enkelt Hjerte, denne Sjæl, \emph{dette} subjektive Hjerte
begjærer, er ligegyldigt. Om et enkelt Hjerte brister af
Sorg --- en enkelt Flue drukner i en Flødepotte --- hvad
Magt ligger vel derpaa? Der er Hjerter nok i Hjertets
Begreb, thi Hjertets Begreb, det almene Hjerte, omfatter
jo alle mulige Hjerter.

I denne Aand taler Hegel, naar han i anden Del
af sin Æsthetik karakteriserer Middelalderens Riddervæsen
og den romantiske Kjærlighed. Subjektivitetens
Valg viser sig for ham som en Egensindighed, en »Partikularitetens
Halsstarrighed«. »Kjærlighedens Kollisioner«,
siger han, »især naar de stilles kæmpende imod substantielle
Interesser, beholde altid et Anstrøg af Tilfældighed og Uberegnelighed;
det er Subjektiviteten som saadan, der med
sine ikke i og for sig gyldige Fordringer sætter sig op
imod det, som for sin egen Væsenligheds Skyld har
% --- p. 594 ---
\setcounter{footnote}{0}%
\opage{594}Fordring paa Anerkjendelse. Individerne i de Gamles
høie Tragedie, Agamemnon, Klytemnestra, Orest, Ødipus,
Kreon o. fl., have vel ligeledes individuelle Formaal,
men det Substantielle, den Pathos, der i Handlingens
Indhold er den bevægende Kraft, har absolut
Berettigelse og er i sig selv af almindelig Interesse\dots{}
Naar Klytemnestras Brøde i dette konkrete Tilfælde ikke
straffes, naar den Krænkelse, Antigone maa lide som
Søster, ikke hævnes, saa er dette en Uret i og for sig.
Men« --- her sigter han til Riddervæsenet --- »disse
Kjærlighedens Lidelser, disse strandede Forhaabninger,
denne Forlibthed overhoved, disse uendelige Smerter,
som en Elskende føler, denne uendelige Lyksalighed og
Salighed, som han forestiller sig, have ingen i sig selv
almindelig Interesse, det er Noget, som kun angaaer
ham selv. Ethvert Menneske har vistnok et Hjerte for
Kjærlighed og Ret til at blive lykkelig ved den, men
naar han her just i dette Tilfælde, under disse og disse
Omstændigheder, netop med Hensyn paa denne Pige,
ikke naaer sit Maal, saa er dermed ingen Uret skeet.
Thi det er ikke i og for sig nødvendigt --- dass er sich
gerade auf dieses Mädchen kapricionire, og saa skulde
vi interessere os for den høieste Tilfældighed, for Subjektivitetens
Vilkaarlighed, der er uden al Udstrækning
og Almindelighed«.

Denne, som det maaske kunde synes, noget Jeronymusagtige
Opfatning af den romantiske Kjærlighed, er
objektiv, strengt videnskabelig. Hvor ægte videnskabelig
den er, kunne vi let bevise. Der skal til Løsning
af en geometrisk Opgave f. Ex. afsættes et Punkt
eller drages en Linie. Men hver Gang et Punkt afsæt%
% --- p. 595 ---
\setcounter{footnote}{0}%
\opage{595}tes, bliver det \emph{dette} Punkt; Linien, der drages, bliver
hver Gang en \emph{denne} Linie. Den, som nu paa Almindelighedens
Vegne har Lyst til at drille Mathematikeren,
kan med en vis Grund gjøre Indvending. Du skal afsætte
et Punkt, det er nødvendigt; men Du afsætter
\emph{dette} Punkt, det er jo tilfældigt; Du vil drage en Linie,
men hvorfor just denne Linie? Du lover os det Almene
og byder os kun det Enkelte, Du sigter paa det
Nødvendige og rammer kun det Tilfældige: det er en
underlig Methode. Mathematikeren vil imidlertid let rede
sig ud af Forlegenheden. Vel maa han indrømme, at
ethvert bestemt Punkt, hvorsomhelst han end afsætter
det, aldrig bliver det almindelige Punkt, og at det Samme
maa siges om Linien; men den Tilfældighed, der saaledes
løber med i Valget af det enkelte Punkt og den enkelte
Linie, er saa langt fra at svække Nødvendigheden,
at den tvertimod selv er nødvendig for at give det Almene
Tilværelse. Da nu det ene Punkt i denne Henseende
kan være ligesaa godt som det andet, den ene
Linie ligesaa tilfredsstillende som den anden: saa kunne
det vilkaarlige, forsaavidt tilfældige, Punkt og den vilkaarlige,
forsaavidt tilfældige, Linie meget vel betragtes
som fuldgyldige Repræsentanter for Punkt og Linie i Almindelighed.

Ifølge Theorien skulde det nu, strængt taget, forholde
sig med Elskeren og hans mulige Valg som med Mathematikeren
og hans mulige Punkter; men det er umuligt.
Den iblandt alle fra Poesien berømte Elskere, der allermest
synes at nærme sig »Væsensuniversalitetens Bestemmelse«
med tilsvarende »Forsigværen i den frigjorte Almindeligheds
Form«, er uden Sammenligning Don Juan.
% --- p. 596 ---
\setcounter{footnote}{0}%
\opage{596}Og dog vil selv en Don Juan med sit frie Valg, sin
næsten uoverskuelige Rigdom af blomstrende Muligheder,
trods den Lethed, hvormed han vælger, og den Hurtighed,
hvormed han vexler, umulig kunne forholde sig til
»Baronesser og Comtesser, Bønderpiger og Prindsesser«
som en Mathematiker forholder sig til Punkter, Flader
og krumme Linier. At han, hver Gang han træder fra Mulighed
ind i Virkelighed, at forstaae, naar han paa Existensens
Vei erotisk vælgende exekverer, maa nøies
med Een ad Gangen, en Enkelt, en Denne, er ikke Vanskeligheden;
thi det Samme kan bogstavelig siges om
Mathematikeren, der maa nøies med dette Punkt, denne
Linie. Don Juan vil indrømme Hegel, at ethvert Menneske
har et Hjerte for Kjærlighed og Ret til at blive
lykkelig ved den, samt at en Elsker ikke for Liv og
Død ubetinget skal binde sig til en Enkelt som til den
Eneste i Verden, eftersom det jo ikke i og for sig er
nødvendigt, »dass er sich gerade auf dieses Mädchen
kapricionire«; og dog er Don Juan ligesaa lidt enig med
Hegel som med Jeronymus: han tilfredsstiller kun det
Subjektive, ikke det Objektive; istedenfor at bringe det
Tilfældige til at gaae op i det Nødvendige, springer han hen
over »Begrebet« saa let som han springer over Kirkemuren;
han parodierer Theorien og bringer hvert Øieblik det Nødvendige
til at gaae op i det Tilfældige, og det uagtet
»Subjektivitetens Vilkaarlighed« netop hos ham har,
hvad Hegel fordrer, »Udstrækning og Almindelighed«.

Det erotiske Valg, hvorved Tilbøielighed og Lidenskab
ofte faae en saadan Overvægt over den kolde Fornuft,
at Hjertet løber af med Forstanden, er jo vistnok
--- siger man --- af saa subjektiv, individuel Natur, at
% --- p. 597 ---
\setcounter{footnote}{0}%
\opage{597}det Personlige her i mange Henseender støder an mod
det rene Begreb, hvisaarsag Theoriens Almene, hvad
dette delikate Kapitel angaaer, helst indskrænker sig til
at være begrændsende uden at være begrundende. Men
efterhaanden som vi komme høiere op i de sædelige Regioner,
vil det Objektive komme til sin Ret; Pligtens
Stemme vil hæve sig og indskærpe »den individuelle Værens
Almindelighed«. Man tænke f. Ex. paa den Pligt
at elske sin Slægt, sit Folk, sit Fædreland. --- Her møde
vi igjen Theoriens gamle Vildfarelse.

Theorien fordrer, at den sande Personlighed skal
bøie sig for Tilværelsens Almene og tage sit Udgangspunkt
i det Objektive, det Nødvendige. Det vil Personligheden
ikke, den protesterer; dens høieste Lov er ikke
objektiv, men subjektiv: den har sit Udgangspunkt i sig
selv. Kun med dette Udgangspunkt, i Kraft af denne Lov,
Selvbestemmelsens, Villiens, Selvudviklingens, Frihedens
Lov, anerkjender den det Objektive og fyldestgjør det
Almene. Tyngdens Lov er objektiv, men Pligtloven er
subjektiv, hin er upersonlig, denne personlig. Det er
Pligt at opgive sin Egensindighed og beherske sit Hjertes
egenkjærlige Tilbøieligheder; det er Pligt at elske sin
Slægt, sit Folk; det er Pligt at offre sig for sit Fædreland.
Men hvor forskjellige Pligterne end ere, saa er dog
Pligtloven een og den samme, den er en indre Personlighedens
Lov, en uforkrænkelig Selvhedslov. Personligheden
giver sig hen; hvorfor? Fordi den i Selvhengivelsen
føler, at den ikke taber, men vinder --- sig selv! Personligheden
opoffrer sig; hvorfor? Fordi den i Opoffrelsens
Øieblik erkjender, at den just nu tilfredsstiller sin
oprindelige Drift, sin inderste, sandeste, dybeste Trang.
% --- p. 598 ---
\setcounter{footnote}{0}%
\opage{598}Naar en Arnold Winkelried i Slagets Trængsel griber
omkring de fjendlige Spyd, samler dem i sit Bryst og
trænger frem: da skeer det ikke af videnskabelig objektive,
nei, af subjektive, ubetinget subjektive, personlige
Grunde. At Helten har indaandet det Objektive og bevæges
af det Almene, idet han bevæges af Ideen, kan
Enhver forstaae; men den bevægende Grund er subjektiv:
den er i Villien, thi den er i Hjertet, der, hvor
ingen Theori naaer hen. Skulde Theorien af objektive
Grunde gjøre det begribeligt, hvorfor just denne Enkelte,
denne Ene, i Farens Yderste offrer sig for alle de Andre:
det vilde blive en vidtløftig Overveielse. Han elsker
Friheden, han elsker Fædrelandet, han elsker de
Kjære i Hjemmet; men Grunden, den objektive Grund?
Det er ligesaa umuligt at angive den almene Grund,
hvorfor en Winkelried saa udelukkende, saa inderlig dybt
er kommen til at elske sit frie Schweitz, som det er
umuligt at udfinde den almene Grund, hvi Romeo kun
kan elske sin Julie, eller Axel sin Valborg.

At de til den saakaldte objektive Ethik hørende
Ideer ligesaa oprindelig røre sig i Sind og Hjerte som
de elskende Individers gjensidige Kjærlighedsfølelse, derpaa
afgiver vor anførte Tragedie jo netop et talende
Exempel. Til at omgaae det objektive Forbud aabner
sig en Udvei; Axels Ven har forberedt Alt til Flugt;
til Afsked knæle de Elskende paa Harald Gilles Grav:
da høres Toner af et Horn langt borte. Axel kjender
disse Toner; det er »en Krigerlur«. Paa Vilhelms Opfordring:
»følg hvorhen din gode Norne vinker«, svarer
han: »min gode Norne vinker mig til Pligt.« Nu begynder
Idebevægelsen personlig fra et nyt Udgangspunkt;
% --- p. 599 ---
\setcounter{footnote}{0}%
\opage{599}men Axel er øieblikkelig enig med sig selv: her er slet
ingen Strid mellem Objektivt og Subjektivt; kun en
indre Forvandling, en ethisk Omsætning af personlige
Grunde. »Nu kan jeg ikke reise; det er umuligt, plat
umuligt, Valborg! Vilhelm. Du vil? Axel. Forsvare
Hakon Herdebred. Vilhelm. Din Avindsmand? Axel.
Hvad Avindmand? Ham har jeg trodsig Ryggen vendt;
nu er min Konge bestædt i Nød, ham skylder jeg mit
Liv. Vilhelm. Du skylder Fædrelandets Vel dit Liv, ei
Hakon. Mener Du at gavne Norge med at forsvare ham,
den unge Daare, som overtræder alle Pligter for sin Lidenskab?
Som ringeagter Dyden og Undersaattens Ret?
Som glemmer ganske det Heles for sin egen Tarv. Nei
Axel! hvo Konge være vil, ham bør det sig at tænke,
handle kongeligt. Kom! Erling er en berømt og tapper
Høvding, ædel. Forsink ham ikke Seiren; han vil ei
misbruge den. Lad Norges Scepter funkle fra Heltens
Haand; det ruster i en Niddings«. --- Det er en egen
Vending, Sagen nu har taget. Saalænge Hindringen laa i
det kanoniske Forbud, kunde det synes, som om Brydningen
mellem det Personlige og Upersonlige fremkom
derved, at Tilbøieligheden stred imod Pligten, som var
det kun subjektiv Lidenskab, der vilde have det Objektive
bort, fordi almene, i sig gyldige Grunde støde fjendlig
op mod Hjertets kjæreste Ønske. Her er nu Forholdet
lige det omvendte; nu er det netop det Objektive,
der er paa Tilbøielighedens Side og begunstiger
Ønsket. I og for sig, vil sige: i objektiv Forstand, er
der ingen tvingende Grunde, hvi Axel skulde blive og
tage Del i Kampen. Hvorfor skulde han vel gjøre det?
Det kan ikke være for Hakons Skyld, thi i og for sig
% --- p. 600 ---
\setcounter{footnote}{0}%
\opage{600}har Hakon store Feil, og Erling store Fortrin. Det kan
ikke være for Fædrelandets Skyld, thi dersom det tør
forudsættes, at Erling muligvis vil blive en bedre Konge
end Hakon, saa har Fædrelandet jo i og for sig ikke
tabt, men vundet, hvis Erling seirer. Men er det muligt,
ja endog sandsynligt, at Axel snarere skader end
gavner Norge ved at blive og kæmpe for Hakon: saa
kan han da heller ikke være forpligtet dertil for Ærens
Skyld, thi at skade sit Fædreland i den Indbildning, at
man gavner det, er i og for sig ingen Ære. Sagen er
i al Fald saa politisk betænkelig, at Vilhelm, i det Mindste
indtil videre, maa synes at have hele Retten paa sin
Side; vi kunne med en lille Forandring anvende hans Ord:
»Naar Alt er lykket, naar tryg den høie Schwarzburg
med sin Ringmur beskytter os; naar Fader Erland har i
Borgkapellet, hvor Leonora tvunget blev til Rudolf, lagt
Eders Hænder i hinanden, bli'er Stunder nok« --- til
Overveielse. Hvor er saa den for Axel sande afgjørende
Grund? Den er i hans eget Selv, hans stærke Æresfølelse,
hans ridderlige Tænkemaade; Ridderen erkjender
sin Pligt: Forpligtelsens Grund er personlig.

Stille vi nu det Patriotiske mod det Erotiske, saaledes
at hint betegner det yderste Objektive, dette det yderste
Subjektive, da maae begge Yderled nødvendigvis udgaae fra
det sædelige Selv; begge have deres Midtpunkt i Sædelighedens
Ide og Sædelighedens Ide i Personligheden.
I skarp Modsætning til Individernes Vilkaarlighed, tilfældige
Lyster og Indfald har Sædelighedsideen tilvisse
sine ufravigelige Pligtlove og i Pligtlovenes almene Gyldighed
sit Objektive; men det maa paa det Eftertrykkeligste
indskærpes, at disse Pligtlove alle uden Und%
% --- p. 601 ---
\setcounter{footnote}{0}%
\opage{601}tagelse ere Selvhedslove, altsaa i strengeste Forstand
Frihedslove, at Forpligtelsen kommer indenfra, fra Hjertets
Grund, fra Villien; thi Villien er i Hjertet, og som
Hjertet er, saaledes er Villien. Kant har med Rette
fremhævet, at Agtelse for Pligtens Bud er uadskillelig
fra Agtelse for sig selv, at Moralloven i Menneskets
Bryst maa ligesom Stjernehimlen over hans Hoved vække
hans høieste Beundring.

Kunde det nu med Rette siges, at Menneskets Selv
er sit eget Princip, at det, hvad den atheistiske Ethik
forudsætter, udelukkende danner sig ved sig selv, udvikler
sig selv, opholder sig selv o. s. v., da vilde unægtelig
heraf følge, at Pligterne mod vor Næste uden
videre gik op i Pligterne mod os selv; og Regnskabet
skulde nok vise, at Enhver er sig selv nærmest. Men
det er netop en Grundmodsigelse at ville hævde det
menneskelige Selv ved Udelukkelse af det guddommelige.
»Det Godes Ide er Villiens Ide: til det absolute
Gode svarer den absolute Villie, det absolute Selv.
Det absolute Selv er den hellige Gud. Ligeoverfor
det absolute Selv staaer saa, med Verden til Mellemværende,
Mennesket, det relative Selv, det Selv, der
ikke er godt, men forpligtet til det Gode, den relative
Villie, der ikke er hellig, men syndig«\footnote[1]{Om Hindringer og Betingelser for det aandelige Liv. S. 101.}. Her
er det Vendepunkt, hvor Sædeligheden viser tilbage til
Religionen som til sit oprindelige Grundlag: Sædelighedsprincipet
staaer og falder med Religionen. At Religionen
maa bære Sædeligheden, kan paavises, men
ikke bevises; her er den Kjendsgjerning, som Videnska%
% --- p. 602 ---
\setcounter{footnote}{0}%
\opage{602}ben kan belyse, men ikke begrunde; og hvorfor ikke?
Fordi Principet er personligt.

Skulde en objektiv Bevisførelse med tilsvarende
Begrundelse være mulig, da maatte den fra Logiken
velbekjendte Modsigelsens Grundsætning afgive Rettesnoren.
Men Enhver, der kjender lidt til Kulturhistorien,
maa kunne sige sig selv, at det ingenlunde forholder
sig med Livsanskuelser som med videnskabelige Theorier:
Aandsbegreberne bøie sig altid let efter den engang
valgte Livsanskuelse, men aldrig omvendt. Der høres Stemmer,
som raabe: »Vi have ingen Plads for en villende
Gud, vi føle ingen Trang til personlig Udødelighed, vi
agte ikke paa subjektive Grunde, vi ynde ikke den omvendte
Videnskab« o. s. v. Det vil da sige med andre Ord,
at der gives tænkende Individer, hvis ethiske Selvfølelse
er saa yderlig svag, at de midt i den objektive Verden
ikke kunne mærke Subjektivitetens Selvstændighed, men
betragte Frihedslovene i Aanden som en theoretisk Afspeiling
af Nødvendighedens Love i Naturen. Herimod
er fra Videns Standpunkt Intet at indvende. Den, der
absolut vil aflede sit Selv af det Selvløse og bøie
sig personlig under det Upersonlige, lad ham bøie sig,
lad ham aflede, det bliver hans egen Sag. At afvise
Religionen bør Enhver i Viden have Lov til; men
naar Nogen i sit Selvbedrag gaaer saa vidt, at han i
Videnskabens Navn vil bilde Verden ind, at han virkelig
kan modbevise Religionens Princip, da maa der
protesteres; thi at afvise er ikke det Samme som at
modbevise.

Den haardeste Modsætning og dog umiddelbareste
Sammenblanding af subjektive og objektive, personlige og
% --- p. 603 ---
\setcounter{footnote}{0}%
\opage{603}upersonlige Grunde finde vi i Stridighederne om de historisk
bestaaende Religioner. En Muhamedaner, en Jøde,
en Kristen ville let blive enige om, at der kun kan være
een eneste sand, saliggjørende Tro, at Religionens Sandhed
er den almene, i og for sig gyldige Sandhed, at
alle Mennesker uden Persons Anseelse ubetinget maae
bøie sig for den ene, af Gud selv aabenbarede og desaarsag
ubedragelige Sandhed: den afgjørende Vægt falder
forsaavidt paa det Almengyldige, det Objektive. Men
naar der saa spørges, hvilken Aabenbaring, hvilken Tro,
der da i og for sig er den ubetinget sande, saa opløser
Enigheden sig i Uenighed. Enhver af de Trende har sin
Religion og sin Tro; den Ene er ligesaa forvisset om
at have Sandheden paa sin Side som den Anden. De
ivrige Religionsbekjendere kunne stride med hverandre;
de kunne i opbrusende Lidenskab myrde hverandre,
men overbevise hverandre af objektive Grunde kunne
de ikke.

Glemmes det i Stridens Hede, at Troens personlige
Sandhed maa begynde med Selvet, det Subjektive, og
tilegnes med Frihed, saa skyder Lidenskaben sig øieblikkelig
ind under det Objektive, det Upersonlige, saa
vil den Ene tvinge den Anden ved objektive Beviser,
med objektiv Magt, og saa have vi \emph{Fanatismen}\footnote[1]{„Troesfanatismen er bekjendt nok; Særkjendet er, at den selvsyge Subjektivitet i lidenskabelig Forblindelse vil forvandle den subjektive Sandhed til en objektiv Magt, en Tilværelseslov, for hvilken alle andre Subjektiviteter ubetinget skulle bøie sig. Vidensfanatismen er mindre bekjendt; men derfor ikke mindre karakteristisk. Den har naturligvis ikke sit Udspring fra Videnskaben, men fra en selvisk Fordreielse af Videnskaben. Istedenfor at tage Erkjendelsen af det Personlige for hvad den er, en subjektiv praktisk Erkjendelse, gjør den Vidende Selverkjendelsen uden videre objektiv theoretisk. Det er et stort psykologisk Selvbedrag. Den Vidende bilder sig ind, at den Subjektivitet, hvorfra han gaaer ud, er et objektivt almengyldigt Væsen, en ubedragelig Methafysikbevidsthed, en sand Normalsubjektivitet, medens den dog i Virkeligheden kun er en Abstraktion af den egne individuelle Subjektivitets tilfældige Tilstand. Langt fra at befries fra „Egoismen“, bliver den Vidende ved dette Selvbedrag tvertimod bedaaret af Egoismen; thi at gjøre den egne Subjektivitets tilfældige Hang til objektiv Maalestok for alle andre Subjektiviteter, er Egoisme. Hvad Normalsubjektiviteten, „i Kraft af Tilværelsens objektive Princip“, har Trang til, det skal nu hver menneskelig Subjektivitet væsenlig have Trang til, nemlig at blive metafysisk; hvad Normalsubjektiviteten ikke føler Trang til, en personlig Gud, en Trøst i Bønnen, et evigt Liv o. s. v., slige Drømmerier og subjektive Kjællingerier tør nu ingen sand --- objektiv sand --- menneskelig Subjektivitet for Fremtiden føle Trang til. At føle Trang til Sligt er jo at afvige fra det Normale, at afvige fra det Normale er at afvige fra det Almene, at afvige fra det Almene er at afvige fra Sandheden: Maalestokken er objektiv. Med sin objektive Maalestok behøver Normalsubjektiviteten nu blot --- i al Upersonlighed at faae et lille Anfald af personlig Sandhedsiver, og den ulmende Vidensfanatisme vil forraade sig ved Angreb paa Troen“. --- Hr. Professor Brøchners filosofiske Kritik gjennemseet af R. N. S. 18---19.}.

% --- p. 604 ---
\setcounter{footnote}{0}%
\opage{604}Fastholdes derimod det Individuelle, det Subjektive med
saa ensidig Frihed, at den personlige Sandheds almengyldige
Indhold falder bort, saa er Forskjellen mellem
de positive Religioner uvæsenlig, og saa have vi Indifferentismen.

Det sextende og syttende Aarhundredes skolastiske
Orthodoxi gav den protestantiske Troesfanatisme et
theologisk Grundlag; at undergrave dette Grundlag var
en Opgave for Rationalismen, men Opgavens Løsning
førte til Indifferentisme. Den aandrigste Repræsentant
for Rationalismen i forrige Aarhundrede er uden Sammenligning
G. E. Lessing.

Det Skrift, hvori Lessing med størst Virkning har
bekæmpet Fanatismen og den orthodoxe Sneversindet%
% --- p. 605 ---
\setcounter{footnote}{0}%
\opage{605}hed, er vistnok hans »\emph{Nathan der Weise}«. Fablen
er henlagt til Korstogenes Tid. Saladin har erobret
Jerusalem; en ung Tempelherre, som er forelsket i Nathans
adopterede Datter Recha, henvender sig til Patriarken
om Raad. Som en Legemliggjørelse af Kirkens
over al Fornuft ophøiede Myndighed fordrer Patriarken,
at hans Raad skal følges i blind Tro, og kommer derfor
i nogen Forlegenhed, da han ikke kan faae at
vide, om det foreliggende Tilfælde: at en Jøde har med
faderlig Ømhed opdraget et kristent Barn i Kjærlighed
til alt Godt, er et Faktum eller blot »et Problem«. Tempelherren,
der kun saadan i Almindelighed vilde høre den
højærværdige Mands Mening, mener, at det kunde være det   % PRINTED AS IS (p. 605): „højærværdige“ with j
Samme. »Det Samme«? udbryder Patriarken; »der seer
Herren, hvor dog den stolte menneskelige Fornuft kan
fare vild i geistlige Sager! --- Paa ingen Maade!» Patriarken
har hverken Tid eller Lyst til at befatte sig
med Tankeproblemer --- Fornuftundersøgelsernes pro et
contra henviser han til Theatret --- han aner et haandgribeligt
Faktum og vil i hellig Iver øieblikkelig have
fat paa Jøden. »Den pavelige og keiserlige Ret bestemmer
jo, at en Jøde, der forfører en Kristen til Frafald,
skal paa Baalet. Saameget mere da en Jøde, der har
øvet Vold mod et kristent Barn«. Men, bemærker Tempelherren,
naar Jøden nu af Barmhjertighed har antaget
sig Barnet, som ellers maatte være omkommet i
Elendighed? --- »Siger Intet! Jøden skal brændes. Bedre,
at det var omkommet i Elendighed, end at det til
sin evige Fordærvelse, blev frelst saaledes. Desuden,
hvad Ret har Jøden til at foregribe Gud? Gud kan,
naar han vil, frelse uden ham«. --- Ogsaa trods ham,
% --- p. 606 ---
\setcounter{footnote}{0}%
\opage{606}skulde jeg mene, gjøre salig. --- »Siger Intet! Jøden
skal brændes.« --- Det gaaer mig til Hjerte, vedbliver
den unge Mand, især da man siger, at Pigen ikke
er opdraget i nogen bestemt Tro, men kun har hørt om
Gud, hvad Fornuften lærer. --- »Siger Intet! Jøden skal
brændes.« Da Tempelherren ikke vil gaae til Bekjendelse
og sige rentud, hvem den Paagjældende er, vil
Patriarken strax til Sultanen; thi Saladin, har ifølge
Kapitulationen, svoret at beskytte den katholske Religion
med alle dens hellige Rettigheder.

I faa, men levende Træk har Lessing her udkastet
et Billede af en ægte Fanatiker; om Fanatikeren tilfældigvis
er Patriark i Jerusalem eller maaske »Hauptpastor
in Hamburg«, gjør jo intet til Sagen, naar kun Tegningen
er korrekt. --- Den blinde Autoritet, den blinde Tro
og den blinde Iver ere saa inderlig forbundne, at man
fra to af disse Mærker altid kan slutte sig til det
tredie. Idet Kirken med sine ufeilbare Dogmer opstilles
som Gjenstand for ubetinget Underkastelse i Tro, svinder
det Subjektive umiddelbart hen i det Objektive,
det Personlige i det Upersonlige. Den, der ikke troer
paa Kirkens Ufeilbarhed, bliver evig fordømt. Nu kan
jo vistnok Angst for en evig Fordømmelse siges at være
en subjektiv Grund; men hvad er det, der vækker denne
Angst? Kirkens myndige Ord, et Ord mere objektivt
end Solen paa Himlen. Saa standser Selvudviklingen ved
det første Skridt; thi Troens første Skridt er, at Mennesket
giver Kirken sin evige Sjæl, sit inderste Selv, sin Samvittighed.
Enhver nærmere Forklaring, enhver Grund, ethvert
bestemmende Hvorfor er at søge i Kirkens Lære, dens
Forjættelser, Bud og Forskrifter: den Troende beholder
% --- p. 607 ---
\setcounter{footnote}{0}%
\opage{607}Intet for sig selv uden den blinde Hengivelse, den ubetingede
Lydighed. Men netop ved denne Selvhengivelse
eller rettere Selvopgivelse bliver den religiøse Subjektivitet
et Organ, et villieløst og dog villiestærkt Redskab
for det i sig Guddommelige, det Objektive. Den Troendes
Sag gjøres nu uden videre til Guds Sag, hans Villie
til Guds Villie, hans Venner til Guds Venner, hans
Fjender til Guds Fjender. Saa tændes Lidenskaben, og
Troesraseriet begynder; thi hvilken Straf kan vel være
for haard, naar det gjælder Guds Fjender? De maae
forfølges med Ild og Sværd, hudstryges, korsfæstes,
brændes: ethvert Middel til deres Udryddelse er Gud
velbehageligt. Patriarken i Jerusalem skulde Intet have
imod, om Sultan Saladin, skjøndt han selv er en Vantro,
der hundede Gange har fortjent Baalet, vilde laane   % PRINTED AS IS (p. 607): „hundede“ for hundrede
den hellige Kirke sin stærke Arm og være den behjælpelig
med at brænde alle Jøder.

Den religiøse Fanatisme er en ubændig Troesegoisme.
I samme Grad som Selvforstaaelsen mangler og de personlige
Grunde tabe sig i blind Underkastelse under en
foregiven, udvortes guddommelig Autoritet, stiger den
selviske Lidenskab. Der gaaer Ild i det naturlige Selv.
Ildfænomenet er Tegn paa en psykisk Forbrænding, hvori
det Upersonlige og det Personlige gjensidig hidse og fortære
hinanden. Den ophidsede Egenvillie bilder sig ind at gjøre
Guds Villie og stride en personlig Strid for den personlige
Gud. Et forfærdeligt Selvbedrag! Fanatikerens Gud,
den objektive og dog grundløse Hersker, er i sin blinde
Vilkaarlighed et selvløst Selv, en ubønhørlig Magt, der
tilintetgjør Friheden, en giftig Aande, der forpester Personligheden.
Dette fanatiske Væsen, dette for Civilisa%
% --- p. 608 ---
\setcounter{footnote}{0}%
\opage{608}tion og Humanitet fjendlige Princip, lige farligt og for den
ædlere Selvfølelse lige oprørende, enten det, saaledes som
paa Korstogenes Tid, slaaer ud i vilde, fortærende Luer,
eller det under Masken af protestantisk statskirkelig
Rettroenhed som en ulmende Glød pustes op ved dogmatiske
Inkvisitioner og giver sig Luft i borgerlige Plagerier,
dette Dogmatismens Uhyre, mod hvilket Lessing
sikkert, behændigt i mangen Dyst havde svunget Kritikens
Klinge, er det, han med æsthetiske Vaaben vil
tillivs i sin \emph{Nathan} den \emph{Vise}.

Samtalen mellem Nathan og Saladin angaaende
Forholdet mellem de tre Religioner, den muhamedanske,
den jødiske og den kristelige, er høist sindrig; vi
tænke paa Lignelsen om de tre Ringe. Nathan fortæller.
I Fordums Dage, i den graa Oldtid, levede
en Mand i Østen, som fra en kjær Haand besad en
Ring af uskatterlig Værdi. Stenen var en Opal, der
spillede i hundrede skjønne Farver, og havde den hemmelige
Kraft, at gjøre hver den, som bar den med fuld
Tillid, velbehagelig for Gud og Mennesker. Hvad Under,
at hin Mand i Østen aldrig lod den komme af sin
Finger, og traf den Foranstaltning, at den for evig
skulde bevares i Familien. Nemlig saaledes: Han bestemte
Ringen for den af sine Sønner, som han elskede
mest; og fastsatte, at denne igjen skulde overlade den
til den af sine Sønner, der uden Hensyn til Fødsel, var
ham den kjæreste. Den Søn, der havde Ringen, var
Familiens Overhoved. Saa kom da denne Ring fra Søn
til Søn, omsider til en Fader, der havde tre Sønner,
som alle tre vare ham lige lydige og følgelig hans Hjerte
lige kjære. Det skete da, at han fra Tid til Tid, efter%
% --- p. 609 ---
\setcounter{footnote}{0}%
\opage{609}som han talede i Enrum med hver af Sønnerne især,
og hver Gang fandt, at den han talede fortrolig med,
stod hans Faderhjerte nærmest, kom til at love Ringen
til dem alle tre. Da han nu paa sit Dødsleie kun
havde een Ring, men tre Sønner, og ikke nænnede at
skuffe Nogen af de Kjære, sendte han hemmelig Bud
til en Konstner og lod ham gjøre to andre Ringe, der
i Et og Alt fuldkomment lignede den ægte.

Saladin, der noget utaalmodig har hørt paa dette
»Eventyr«, studser, da Nathan erklærer ham, at den rette
Tro er ligesaa ukjendelig, som den ægte Ring. --- Hvad,
siger Sultanen, skal det være Svar paa mit Spørgsmaal?
Driv ikke Gjæk med mig; jeg skulde dog mene, at de
Religioner, jeg har nævnet Dig, vel ere til at adskille,
at Forskjellen endogsaa gaaer indtil Klædedragt, Spise
og Drikke. --- Men, svarer Nathan, hvad Grundene angaaer,
er der ingen Forskjel. Thi grunde de sig ikke
alle paa Historie, mundlig eller skreven? Men Historie
maa dog vel antages paa Tro og Love? Hvis Ord og
Tro drager man da mindst i Tvivl? Dog vel sine Egnes?
dog vel deres, hvis Blod vi ere? deres, som fra vor
Barndom af have givet os Prøver paa deres Kjærlighed?
som aldrig have skuffet os, uden i Tilfælde, hvori det
var os tjenligt at blive skuffede? Hvor kan jeg troe
mine Fædre mindre end Du dine? Eller omvendt. Kan
jeg forlange af Dig, at Du skal erklære dine Forfædre
for Løgnere, for ikke at modsige mine? Eller omvendt.
Det Samme gjælder om de Kristne. Ikke sandt? ---
Denne Argumentation gjør et stærkt Indtryk paa Saladin;
han udbryder: Ved den levende Gud, Manden
har Ret. Jeg maa forstumme.

% --- p. 610 ---
\setcounter{footnote}{0}%
\opage{610}Spørgsmaalet om Religionens objektive Grundlag
falder altsaa sammen med Spørgsmaalet om historiske
Kjendsgjerningers Troværdighed og da særlig om Forholdet
mellem religiøs og historisk Tro. I en lille Afhandling:
\emph{Über den Beweis des Geistes und der}
\emph{Kraft}, havde Lessing allerede opstillet den for hans
Synsmaade saa karakteristiske Sætning, at tilfældige
Historiesandheder aldrig kunne blive Bevis for nødvendige
Fornuftsandheder. Det er nu mærkeligt at see
Sammenblandingen af upersonlige og personlige Grunde.

De historiske Begivenheder ere udvortes objektive
og tilhøre det Forbigangne; deres Troværdighed kan,
især naar Talen er om Fænomener af det Overnaturlige,
om Tegn og underlige Gjerninger, let blive tvivlsom.
Det Høieste, hvortil man ad denne Vei kan naae,
er kun en vis Grad af Sandsynlighed. Men den blotte
Sandsynlighed giver aldrig fuld Vished; hertil kommer,
at Underet, det fra Naturens Orden Afvigende, altid
maa vise sig som det Usandsynlige. Dommen over
det Objektive bliver paa denne Maade afhængig af Følelse,
af Stemning, altsaa af Subjektiviteten. Og naar saa
den paa sig selv usikkre Subjektivitet griber efter det Objektive,
hvad er da naturligere, end at Enhver, som Nathan
siger, holder sig til det fra Fædrene Overleverede og helst
troer sine Egne? Det historisk Objektive er, at man
fødes, om just ikke med, saa dog til en bestemt Religion:
den fødte Muhamedaner troer paa Muhamed, den
fødte Jøde paa Jehovas Lov, den fødte Kristne paa Kristus.

Paa disse Vilkaar have Individer af forskjellig
Troesbekjendelse ikke Noget at lade hinanden høre. Er
det ikke muligt af historisk objektive Grunde at bevise,
% --- p. 611 ---
\setcounter{footnote}{0}%
\opage{611}hvilken Religion der i og for sig er den sande, saa kan
den ene jo forsaavidt være ligesaa god som den anden.
Det Objektive har givet efter og bøiet sig ind under det
Subjektive. Religionen er bleven et Pietetsanliggende,
en Hjertesag, og Valget er frit. Istedenfor at lægge Vægt
paa Dogmerne, lægger man nu udelukkende Vægt paa
det fromme Sindelag, paa Dyd og gode Gjerninger.
»Velan«! saa taler Dommeren i Nathans Fortælling til
de tre Brødre, der kives om den ægte Ring --- »saa stræbe
da Enhver at lægge sin ubestikkelige, fordomsfrie
Kjærlighed for Dagen, at Stenen i hans Ring maa vise
sin Kraft; han komme denne Kraft til Hjelp med Sagtmodighed,
med hjertelig Fordragelighed, med Velgjørenhed,
med inderlig Hengivenhed i Gud: jeg indbyder da
Eder at møde her igjen, naar Tusinder af Aar ere
svundne; en Mand, der er visere end jeg, vil da sidde
paa denne Stol og afsige Dommen«. Saladin, fortsætter
Nathan, dersom Du føler Dig at være denne vise Mand,
den Forjættede\dots{} Denne Henvendelse gjør et saa mægtigt
Indtryk paa »de Troendes Behersker«, at han styrter
hen til Jøden, griber hans Haand og udbryder bevæget:
»Jeg Støv? Jeg Intet? O Gud«! --- Med en saadan
Selvbekjendelse er Fordragelighedens Lov underskreven;
Fanatismens Baal slukkes: Humaniteten har seiret.

Men ved denne Sagens Vending ere vi komne til
en modsat Yderlighed: til Erstatning for Fanatismen
have vi nu faaet Indifferentismen og med denne Tvivlesygen,
Skepticismen. Er det først almindelig indrømmet,
at ingen af de positive Religioner, der udgive sig
for aabenbarede, er i Stand til at føre et historisk uomstødeligt
Bevis for sin Aabenbarings Ægthed: da følger
% --- p. 612 ---
\setcounter{footnote}{0}%
\opage{612}ganske rigtig heraf, at Bekjendere fra de forskjellige Troessamfund
maae gjøre hverandre Indrømmelser. En Bekjender,
f. Ex. af den kristne Religion, maa da sige til sig
selv: Grunden, hvorfor jeg er Kristen, er, naar jeg skal
være oprigtig, egenlig den, at jeg er født i et kristent
Samfund af kristne Forældre; var jeg bleven født i et
jødisk Samfund, var jeg naturligvis bleven Jøde, og
hvis jeg var bleven født midt i Tyrkiet, ja saa var jeg
bleven Tyrk. Af en saadan Tilstaaelse følger ikke blot
Overbærenhed mod anderledes Troende, altsaa Tolerance,
men ogsaa Ligegyldighed mod de positive Religioners
indbyrdes Forskjellighed.

Resultatet af det forrige Aarhundredes aandelige
Kampe var et afgjørende Brud med den blinde Autoritet,
en Emancipation fra det i Historien udvortes Overleverede.
Man skulde nu mene, at en saadan Emancipation
fra det Udvortes maatte have ført til dyb, indtrængende
Erkjendelse af det Indvortes, det Personlige;
men det skete ikke. Der skete noget Andet. Just
som den fra sit gamle dogmatiske Fangebur bortflagrende
Subjektivitet begyndte at svæve høit over alle faste
Lærebygninger og i Nydelsen af den friske Luft med
den frie Udsigt gav sig hen til allehaande religiøse Forestillinger,
stolende paa, at Enhver blev salig ved sin Tro,
gjorde Fornuften sin guddommelige Ret gjældende og
satte paany det Objektive i Kraft, om end under en
ganske egen Form. Vi skulle her ikke opholde os ved
en udførlig Karakteristik af den saakaldte Fornuftreligion,
men blot fremhæve et eneste Punkt for at vise,
hvorledes Rationalismen, nær ved at ramme det Rette,
skjød Sandheden forbi, idet den sigtende paa Friheden
% --- p. 613 ---
\setcounter{footnote}{0}%
\opage{613}kun traf Nødvendigheden, og styrende mod det Personlige
strandede paa det Upersonlige\footnote[1]{Jvfr. Religionsfilosofi S. 28—35; S. 88---97.}. Grundpunktet, hvori det Objektive mødes med det
Subjektive, er Selvet. Det menneskelige Selv bærer i
sig et dobbelt Princip, et dobbelt Anlæg, et Anlæg for
Viden og et Anlæg for Tro. Naar da den opvaagnende
Fornuft sætter Selvtænkningen i Bevægelse, kommer
det an paa, om Blikket oprindelig er rettet mod
det Objektive eller mod det Subjektive.

Er Blikket rettet mod det Objektive, saa blive ikke
blot Tingene i Verden, men ogsaa Sjælens Indre, Selvets
Væsen, Gjenstand for Viden. Ingen, der taler
naturlig, vil sige: Jeg troer, at jeg er til, at jeg føler,
seer, hører, tænker o. s. v.; han vil med umiddelbar
Forvisning sige: jeg veed det. Det Samme gjælder,
naar Talen er om Andres Bevidsthed. Ingen vil sige:
Jeg troer, at disse Mennesker, som her omgive mig,
virkelig ere levende, at ogsaa de føle, see, høre, tænke
ligesaa vel som jeg; men han vil, forudsat, at han
ikke har en metafysisk Skrue løs, sige ligefrem: Jeg
veed det. Der er et System af Almenbestemmelser,
der ligesaa nødvendig høre til Selvbevidsthedens Begreb.
som Sider og Vinkler høre til Trekanten.

Anlægges nu Begrebsudviklingen saaledes, at Spørgsmaalet
slet ikke er om dine og mine Følelser, dine og
mine Tanker, din og min Villie, dit og mit Selv, men
om menneskelige Følelser og Tanker i al Almindelighed, om
den almene Villie, det almene Selv, saa antage de psykologiske
Begreber et objektivt, videnskabeligt Præg,
% --- p. 614 ---
\setcounter{footnote}{0}%
\opage{614}saa skinner Fornuftens Lys i Bevidstheden. I denne
Belysning ville da ogsaa Sædelighedens Love tage sig
ud som rene Nødvendighedslove; forsaavidt her er Tale
om Tro paa Gud som sædelig Lovgiver, vil denne Fornufttro
være af en saa almengyldig Beskaffenhed, at
det livagtig seer ud, som kunde dens Indhold med næsten
mathematisk Nøiagtighed begrundes og demonstreres af
Forstanden.

Rationalismen har ganske rigtig indseet, at det,
hvorpaa det ved al Religion og Sædelighed væsenlig
kommer an, er en klar Erkjendelse af det menneskelige
Selv, dets Udspring fra Gud, dets evige Bestemmelse,
dets uendelige Perfektibilitet; Rationalismen har
indseet, at den sande Tro maa være noget Andet end
simpel Antagen af det Ubegribelige paa blind Autoritet;
den har indseet, at Religionens Sandheder bør, ligesom
andre Sandheder, prøves efter Selvtænkningens
Princip og bekræftes ved Indsigt, ved Overbevisning.
Men Grundvildfarelsen er iøinefaldende.

Rationalismens Feilgreb er, at den tager Skyggen
for Virkeligheden, Selvets Tankevæsen for det levende
Selv i Kjød og Blod. Ved dette Feilgreb bliver al Selvtænkning
til videnskabelig Tænkning, og Tænkningens
Love falde sammen med den naturlige Verdens Love.
Hvormeget de gamle Rationalister end ivrede for at
holde fast ved Gud, ved Dyd og Udødelighed, det
hjalp ikke. Gud, Dyd og Udødelighed ere Personlighedskategorier,
der vise tilbage til Subjektivitetens
Selvstændighed; men den Subjektivitetens Selvstændighed,
hvorefter Rationalismen greb, skylledes bort af
den objektive Videns Strøm: kun Nødvendigheden selv,
% --- p. 615 ---
\setcounter{footnote}{0}%
\opage{615}det Eviges Form, det kolde Almene, blev tilbage. Det
er med Hensyn paa Aandslæren ikke Tilfældighed, men
videnskabelig Følgestrænghed, at Rationalismen opløste
sig i spekulativ Filosofi, og den spekulative Filosofi i
Atheisme.

Skal Atheismen overvindes, maa det skee ved Selvet;
men saa maa Selvet opfattes, ikke som et almindeligt
abstrakt Væsen, men som en levende Existens,
ikke som et Vidensbegreb, men, hvad Personligheden
angaaer, som et Troesbegreb.

At opfatte Selvet som Vidensbegreb er, vi gjentage
det, at see bort fra sin egen tilværende Bevidsthed,
sine egne Følelser, sit eget Hjertes Tanker og gribe
Bevidsthedens Væsen i dets rene Almindelighed, altsaa
upersonligt. Fra Bevidsthedens Væsen i dets Almindelighed
er det Tanken umuligt at naae hen til din og
min virkelige Bevidsthed, umuligt at naae fra det Upersonlige
til det Personlige.

Anderledes, naar det menneskelige Væsen opfattes
fra Troens Synspunkt. Her begynder Enhver med sig
selv. Opgaven er at erkjende sin oprindelige Bestemmelse,
opdage det Evige i sit eget Selv og forvisse sig
om Selvets Evighed: Opgaven er at begynde personligt.
At dette er en Troesopgave, kan let indsees.

Er Menneskets Sjæl dødelig eller udødelig? For at
faae dette Spørgsmaal besvaret, maae vi allerførst gjøre
os Rede for, hvor det hører hen, af hvad Beskaffenhed
det er, og hvad det er for Grunde, af hvis Gyldighed
Besvarelsen afhænger. Spørgsmaalet er personligt; det
vil sige: her kan ikke begyndes med Sjælens Væsen i
Almindelighed; en saadan Begyndelse er falsk, thi den
% --- p. 616 ---
\setcounter{footnote}{0}%
\opage{616}er upersonlig. Af en falsk Begyndelse udledes aldrig
en sand Slutning; en upersonlig Forudsætning kan umulig
afgive en personlig Følgesætning. For at begynde
personlig, maa Enhver begynde med sin egen Bevidsthed,
og spørge sig selv: er min Sjæl udødelig? Tre
Svar ere mulige; den Spørgende kan svare sig selv
\emph{bekræftende}, han kan svare \emph{benægtende} og han kan
forholde sig \emph{tvivlende}. At de tre Svar, saasandt de
skulle have personlig Betydning, maae udgaae fra tre Individer,
tre forskjellige Mennesker, trænger ikke til Bevis;
thi den, der i det ene Øieblik antager, at Sjælen
er dødelig, i det andet, at den er udødelig, er jo usikker
paa sig selv, altsaa tvivlende. Det er endvidere
klart, at medens der ingen nødvendig Overgang kan
gjøres fra det Almene til det Individuelle, saa kan der
derimod gjøres en fri, personlig Overgang fra det Individuelle
til det Almene. Thi dette staaer fast: saaledes
som det enkelte, virkelige Menneske opfatter sit
oprindelige Selv, sit eget menneskelige Væsen, saaledes
maa han nødvendig opfatte Menneskeslægtens Væsen;
han gjør ikke sin tilfældige Individualitet til Maalestok
for andre Individer, det er Selvets Væsen han gjør til
Maalestok; men Selvets Væsen erkjender han i sit eget
Selv: Lysstraalerne fra det Almene bryde sig i Individualitetens
Prisme, samle sig i Personlighedens Brændpunkt.

Vende vi os til de to Individer med de to afgjørende
Svar og prøve, hvorvidt den Ene ved objektive
Grunde er i Stand til at overbevise den Anden, da komme
vi strax til et negativt Resultat. Den, der føler sig absolut
forvisset om Sjælens Udødelighed, med andre Ord,
om Selvets Evighed, har Grunden til denne Forvisning
% --- p. 617 ---
\setcounter{footnote}{0}%
\opage{617}i sig selv og kan umulig hente den andensteds fra.
Hvad der har paavirket ham udenfra, hvilke Tilskikkelser
han har prøvet, hvad han i sig selv har oplevet,
hvilken Strid, han har stridt, hvilke Kriser han har
gjennemgaaet, inden han opnaaede Forvisningen, derom
er Spørgsmaalet ikke; Spørgsmaalet er om den opnaaede
Forvisnings eiendommelige Art og Natur. Naar dette
Synspunkt fastholdes, da er det uimodsigeligt, at Visheden
er Selvvished, en Vished om Selvet, som kun Selvet
giver, og alle tilsvarende Grunde personlige.

Men den, der i Modsætning hertil holder sig forvisset
om, at Sjælen er dødelig, at den dannes med
Legemet, udvikles med Legemet, opløses og forgaaer
med Legemet, at det menneskelige Selv, om det end
er Bærer og Gjennemgangspunkt for evige Ideer, ikke
selv er evigt, skulde han maaske ikke kunne bygge paa
Videnskaben og retfærdiggjøre sin Livsanskuelse af objektive
Grunde? Umuligt! Alle hans Beviser, hvor skarpsindig
han end skiller dem, og hvor dybsindig han end
samler dem, forene sig til et eneste stort Bevis, Beviset
for, at han befinder sig i et stort, psykologisk
Selvbedrag. Bedraget viser sig ikke ligefrem deri, at
han antager Sjælen for dødelig; nei, psykologisk seet,
bestaaer det oprindelige Selvbedrag deri, at den Bedragne
bilder sig ind at kunne retfærdiggjøre sin Anskuelse
videnskabeligt.

Den, der skal dømme om Naturens Love, maa see
bort fra sit eget Selv, for at dømme videnskabeligt;
den, der dømmer om Trekanten, maa see paa Trekanten,
ikke paa Selvet, thi Trekanten er jo ikke Selvet; den,
der dømmer om Grundstofferne, maa see paa Grundstof%
% --- p. 618 ---
\setcounter{footnote}{0}%
\opage{618}ferne, ikke paa Selvet, thi Selvet er intet Grundstof;
ligesaa, naar der dømmes om Jordudviklingen, om Forverdenens
Fænomener o. s. v.: alle slige Domme ere
objektive derved, at den Dømmende seer bort fra sig selv
for at see paa Gjenstanden.

Men naar Dommens Gjenstand er den Dømmendes
eget Selv, saa maa Synspunktet forandres.

Dommen om en Gjenstand maa altid rette sig efter
Gjenstandens Beskaffenhed og ethvert uvedkommende
Hensyn holdes borte. For at undersøge et Grundstof
som f. Ex. Svovlet maa Kemikeren jo vistnok bringe
det ind i mangfoldig Vexelvirkning med andre Stoffer
og betragte det i mange Forbindelser, men med alt dette
er det dog Svovlets Eiendommelighed, dets fysiske
og kemiske Egenskaber, han uafbrudt har for Øie.
Sætte vi nu Selvet istedenfor Svovlet, saa er det jo
dog indlysende, at den, der vil dømme om Selvet,
uafbrudt maa have sit Øie fæstet, ikke paa det rene
metafysiske Selv --- thi naar der spørges om Existens, er
det metafysiske Selv et Fantom, --- men paa det virkelige
Selv, dets Kræfter og Væsen, dets aandige Krav.   % PRINTED AS IS (p. 618): „aandige“
Men det eneste virkelige Selv, Betragteren virkelig kan
undersøge, er uden Modsigelse hans eget. Det virkelige
Selv, der er Dommens Objekt, er ogsaa dømmende
Subjekt.

Selvet dømmer sig selv enten til Livet eller til
Døden! saaledes er det psykologiske Bedrag hævet;
Dommen er subjektiv: Grundene ere personlige.

Rationalismens og den spekulative Filosofis psykologiske
Feilsyn ligger nu klart for Dagen. Man vilde,
ganske rigtig, begynde med Selvtænkning og gaae til%
% --- p. 619 ---
\setcounter{footnote}{0}%
\opage{619}bage til Selvet; men forvexlende Selvets Almenbegreb
med det virkelige Selv, greb man Viden og mente
Troen, og Følgen blev, at Sjælens Udødelighed svandt
hen i Fornuftens Evighed, og Personlighedens Ret gik
tabt i det Upersonlige.

De, der erkjende, og de, der fornægte Sjælens
Udødelighed, kunne ikke overbevise hverandre psykologisk,
eftersom deres Selverkjendelse er saa grundforskjellig,
at de med Hensyn paa Udødelighedsspørgsmaalet
maae siges at have hver sin Psykologi; men Modsætningen
bliver endnu langt skarpere, naar Talen er
om Religion og Sædelighed.

For i sædelig Forstand at kjende sig selv, maa
Mennesket ville sig selv; for at ville sig selv, maa
han vælge sig selv. Dommen om Selvets Realitet er
i strengeste Forstand en Villiesdom. Den, der erkjender
sin Sjæls Udødelighed, har valgt sig selv i evig
Forstand, har ved dette Valg ikke blot fundet det Evige
i Selvet, men endog, bogstavelig talt, opdaget Selvets
Evighed. Omvendt har den, der erklærer Sjælen for
dødelig, forudsat at han gjør det med fuld Forvisning,
opnaaet denne Forvisning ved en indre, oprindelig Selvbestemmelse, en Villiesakt, et Valg. Saa staaer kun
Tvivleren tilbage; men Tvivleren, der famler mellem
objektive og subjektive Grunde, er ingen Personlighed;
det er en løs, forsvækket Subjektivitet, der løber
vild i sig selv.

Dersom man nu blot vil indrømme os, at Livet i
sig selv ikke alene er noget forskjelligt, men absolut forskjelligt
fra Døden, saa bliver man vel ogsaa nødt til at
indrømme, at den Dom, der gjør Selvet udødeligt, og   % PRINTER'S DEFECT (p. 619): „indrømme“: the i did not ink (a speck); letter supplied
% --- p. 620 ---
\setcounter{footnote}{0}%
\opage{620}den Dom, der gjør det dødeligt, ere absolut forskjellige
Domme, og at Forskjelligheden beroer derpaa, at de
dømmende Subjektiviteter have givet sig ind under absolut
uensartede Principer.

Alle Erkjendelser, der angaae Viden og Videnskab,
høre uden Undtagelse ind under det Princip, der omsætter
det existerende Selv i Tænkningens almene Væsen,
og som vi derfor kalde \emph{Fornuftprincip}; alle Erkjendelser,
som angaae Religion og Sædelighed, høre derimod
ind under det Princip, der sætter det Almene i det
existerende Selv og gjør det levende Menneskes egen
Sjæl til ufravigeligt Middelpunkt, og som vi derfor, naar
der spørges om evigt Liv, kalde \emph{Troesprincip}.

At disse to Principer ikke blot ere uensartede, men
absolut uensartede, kan sees af de tilsvarende Grunde;
thi alle Fornuftgrunde ere i deres almene Gyldighed
objektive, upersonlige, og gjælde i Kraft af Nødvendighed;
men Troesgrundene ere i deres individuelle Gyldighed
subjektive, personlige og svare til Friheden. Den,
der endnu vil mene, at der kun er en relativ, ikke en
absolut Forskjel mellem de to Principer, lad ham forsøge
at aflede Friheden af Nødvendigheden, det Personlige af det
Upersonlige, og han skal visselig strande med Forsøget\footnote[1]{Skal man sammenflikke sig en personlig Livsanskuelse ved at laane Lidt fra Troen og Lidt fra Viden, Lidt fra det Objektive og Lidt fra det Subjektive, Lidt fra det Upersonlige og Lidt fra det Personlige: da bliver det Livsanskuelse derefter.}.

Troen paa Selvets Evighed forudsætter Tro paa
Guds Tilværelse; thi dersom det menneskelige Selv ikke
havde sin Oprindelse fra det guddommelige; dersom
Evighed, Almagt, Alvidenhed, Hellighed, Salighed ikke
% --- p. 621 ---
\setcounter{footnote}{0}%
\opage{621}tilkom det Væsen, der er Altilværelsens Ophav; dersom
den menneskelige Sjæl kun var en organisk Naturfrembringelse:
hvorledes skulde da Sjælen modstaae Forkrænkeligheden,
og Selvet overleve Døden? Men Overbevisningen
om en personlig Guds Tilværelse er i allerdybeste
Forstand personlig.

Saasandt nu Personlighedens Sag, saalænge der er
Rest af Alvor tilbage i Slægten, aldrig kan blive en ligegyldig
Sag, saasandt maa Personlighedsprincipet siges
at indeholde Modgift mod religiøs Indifferentisme; thi
Enhver maa arbeide for Sit, og Regnskabet gjælder
Samvittigheden. Men hvormange Brydninger der end
kunne opstaae mellem Individer med forskjellig Selverkjendelse
paa forskjelligt Udviklingstrin, idet de hver
paa sin Vis ivre for personlige Opgavers Løsning, saa
vil det dog for Fremtiden ikke kunne glemmes, at Personlighedens
Lov er Frihedens Lov, dens Liv et Frihedens
Liv, og Principet selv vil i denne Henseende
indeholde Modgift mod religiøs Fanatisme.

Sand Tro og sand Personlighed ere uadskillelige.
Alt hvad der angaaer Troen, angaaer Personligheden;
hvad der ikke hører med til Begrundelsen af det personlige
Liv, har Intet med Troen at gjøre. Ved at
fastholde denne Grundsætning ville Bekjendere fra de
forskjellige Religioner arbeide sig frem til een og samme
Troesbekjendelse; thi de personlige Livstrin ere forskjellige,
men Personlighedens evige Liv er i Væsen det
samme.

Ophøiet over alle sjæleformørkende Stridigheder hæver
sig i Troens Lys Sandhedens Aabenbarer, Kjærlighedens
Forklarede, Personlighedens evige Forbillede. At kjende
% --- p. 622 ---
\setcounter{footnote}{0}%
\opage{622}Selvets Grund ved at kjende dette Forbillede, at troe
paa Selvets Evighed ved at troe paa dette Forbillede er,
efter Lignelsen hos Nathan, at have den ægte Ring. Naar
saa Dommeren kommer --- om han end tøver i tusinde
Aar --- da vil han vist ikke dømme efter udvortes Tegn,
med blind Autoritet efter objektive Dogmer; men han
vil »aabenbare Hjerternes Raad« og dømme saaledes, at
Enhver under Dommen maa dømme sig selv og give
ham Ret: han vil dømme personlig, i Kraft af personlige
Grunde.

\signature{R. \emph{Nielsen}.}

\end{document}
